% !TeX root = ../Tempel-Daniel.tex
\chapter{Summary}
\label{chap:summary}

This report examined the architecture of modern AI systems through the recurring example of an AI Software Engineering Assistant capable of analysing a repository, modifying source code, executing tests and verifying the resulting implementation. The investigation progressed from the properties of individual language models to the additional components required to operate such models as reliable software systems.

Local LLM deployment is constrained by the relationship between model size, numerical representation and available hardware. Quantization can substantially reduce the memory required for model weights, making larger models practical on limited hardware, but this introduces trade-offs in model quality and inference behaviour. Runtime memory is also affected by the active context and the KV cache, meaning that a model which fits into memory as a file may still require considerably more resources during execution. Model specialisation is therefore relevant in addition to model size: a smaller coding-specialised, instruction-tuned model may be more useful for a software-engineering workload than a larger general-purpose model.

Selecting such a model cannot be reduced to choosing the highest score from a public leaderboard. Benchmarks measure different capabilities under different conditions, and repository-level software engineering introduces requirements that are not captured by general reasoning or code-generation tests alone. Useful model selection should combine published benchmark evidence with controlled evaluation under the intended assistant architecture. Task success, execution time and resource consumption are more meaningful for this workload than an isolated benchmark score.

A capable model alone is insufficient for an operational assistant. Structured outputs and function calling provide machine-readable interfaces through which a model can request actions, while the surrounding runtime remains responsible for validating and executing those actions. This separation is important because structural correctness does not guarantee that an action is semantically appropriate. The model proposes operations, whereas conventional software components retain control over their execution.

The usefulness of each model call also depends on the information available in its active context. Context engineering therefore becomes a separate architectural responsibility. System instructions, task state, interaction history, tool results and persistent memory should not automatically be included in every request. Instead, the runtime should select information relevant to the current decision, preserve important state and compact long-running interaction histories when necessary. Persistent memory, context compaction and prompt caching address different problems and should not be treated as interchangeable mechanisms.

Retrieval extends this approach when relevant information exists outside the active context. For software repositories, retrieval may involve lexical, dense or structural search rather than only a vector database. Repository-aware retrieval must also account for changing working-tree state and distinguish retrieved evidence from the context eventually presented to the model. More complex tasks can benefit from agentic retrieval, in which later search decisions depend on evidence obtained during earlier steps.

The Model Context Protocol provides a standardised integration boundary for exposing tools and resources to AI applications. It can reduce duplicated integration work across heterogeneous systems, but it does not replace backend-specific implementations, define retrieval strategies or provide a security boundary by itself. The Host remains responsible for selecting capabilities, mapping them to the model interface and maintaining a consistent task state across connected systems.

These components can be coordinated through an agent harness. The harness manages repeated model calls, context construction, tool execution and task state, while the model can adapt subsequent actions according to observations obtained during execution. This creates an agent loop rather than a single model response. However, termination of such a loop must remain separate from successful completion: a task should only be considered successful when its acceptance criteria have been verified.

Production deployment introduces further requirements. Model routing can select between eligible local and remote deployments according to capability, latency, cost and data-handling constraints. Efficient model serving depends on mechanisms such as batching, KV-cache management and prefix caching, while long-running agent tasks require durable state, recoverable execution artefacts and bounded retry behaviour. Scaling should follow the actual system bottleneck rather than individual infrastructure metrics in isolation.

Greater autonomy also increases the importance of security. A software-engineering agent may execute repository-controlled code, modify files and interact with external services. Least-privilege tool interfaces, sandboxed execution, network restrictions, credential mediation and approval gates can limit the consequences of incorrect model decisions or indirect prompt injection. Security should therefore remain effective even when the model itself follows a malicious or incorrect instruction.

Finally, reliable operation requires evaluation and observability at the level of the complete system. Model benchmarks alone cannot determine whether an agent successfully completes a multi-step software-engineering task. Evaluation should distinguish failures in retrieval, context construction, tool execution, model decisions and verification. Tracing, regression evaluations, task-success metrics, latency and resource measurements provide complementary evidence about system behaviour.

The central conclusion of the report is that the language model is only one component of a modern AI system. Practical capability emerges from the interaction between the model, context, retrieval, tools, integration protocols, orchestration, infrastructure, security controls and evaluation mechanisms. Building reliable AI applications therefore requires conventional software-engineering principles around explicit interfaces, controlled execution, state management, isolation, observability and verification in addition to selecting a capable LLM.
